% !TeX program = xelatex
% Standalone article. Compile twice; no external bibliography or images.
\documentclass[11pt,a4paper]{article}
\usepackage[margin=25mm]{geometry}
\usepackage{amsmath,amssymb,amsthm,mathtools,esint,mathrsfs}
\usepackage{booktabs,longtable,array,xurl}
\usepackage[unicode,colorlinks=true,linkcolor=blue!50!black,citecolor=blue!50!black,urlcolor=blue!50!black]{hyperref}
\usepackage{xcolor}
\newtheorem{theorem}{Theorem}[section]
\newtheorem{lemma}[theorem]{Lemma}
\newtheorem{proposition}[theorem]{Proposition}
\newtheorem{corollary}[theorem]{Corollary}
\theoremstyle{definition}
\newtheorem{problem}[theorem]{Open problem}
\newtheorem{definition}[theorem]{Definition}
\theoremstyle{remark}
\newtheorem{remark}[theorem]{Remark}
\DeclareMathOperator{\Ric}{Ric}
\DeclareMathOperator{\Rm}{Rm}
\DeclareMathOperator{\Vol}{Vol}
\DeclareMathOperator{\tr}{tr}
\DeclareMathOperator{\Id}{Id}
\DeclareMathOperator{\diver}{div}
\DeclareMathOperator{\dist}{dist}
\DeclareMathOperator{\Orb}{Orb}
\DeclareMathOperator{\Met}{Met}
\DeclareMathOperator{\Hull}{Hull}
\DeclareMathOperator{\Gr}{Gr}
\newcommand{\R}{\mathbb R}
\newcommand{\CP}{\mathbb{CP}}
\newcommand{\cE}{\mathcal E}
\newcommand{\dd}{\,d\mu_g}
\newcommand{\avg}[1]{\fint_M #1\dd}
\newcommand{\norm}[1]{\left\|#1\right\|}
\newcommand{\HS}{\mathrm{HS}}
\newcommand{\op}{\mathrm{op}}
\newcommand{\roundref}[2]{\href{run:../references/rounds/#1}{#2}}
\newcommand{\baseref}[2]{\href{run:#1}{#2}}
\newcolumntype{L}[1]{>{\raggedright\arraybackslash}p{#1}}
\numberwithin{equation}{section}
\setlength{\emergencystretch}{3em}
\allowdisplaybreaks
\hypersetup{pdftitle={Integral Ricci Rigidity: Uniform Obstructions and Extensions},pdfauthor={Qiushi Yau26 research}}
\title{Integral Ricci Rigidity:\\ Uniform Obstructions and Extensions}
\author{Qiushi Yau26 research}
\date{21 September 2026}
\begin{document}
\maketitle
\begin{abstract}
We develop the integral near-equality obstruction obtained from actual
rescaled Seiberg--Witten monopoles. A bounded spinor admits an intrinsic
second-derivative estimate; its quadratic form produces a smooth Weyl
model whose divergence can be used to test the weak Bianchi identity.
The resulting upper estimate contradicts a positive weighted Hodge
energy near the critical Weyl orbit. The argument yields a uniform
all-sign Ricci-energy gap for special monopole classes, with blow-ups
and suitable connected sums as applications. We include the class-square
cost, scalar truncation, topology transfer, and compatibility calculations
needed to use the result. We then formulate further questions concerning
economical class selection, boundary and source defects, multiple
calibrating fields, positive Ricci pinching, and higher dimensions.
These questions remain open. Throughout, the objective is the full
dimension-dependent Einstein-existence question, whose two sign branches
and unrestricted output metric must be kept distinct from the range of
this particular obstruction.
\end{abstract}
\tableofcontents

\section{The underlying problem and its quantifiers}\label{sec:problem}
\subsection{The full objective}\label{sub:full-objective}
The working formulation of Yau26 in this project is the following.
For each integer $n\ge2$, do there exist real constants $c_n<C_n$,
depending only on $n$, such that every closed connected smooth
$n$-manifold carrying a Riemannian metric $g$ with
\[
                 c_ng\le\Ric_g\le C_ng
\]
admits a smooth positive-definite Einstein metric $h$?
The metric $h$ is unrestricted: it need not be close to $g$, have the
same Ricci sign, or preserve a complex structure, fibration, or symmetry.
We denote this statement by $Q_n$. The exact historical correspondence
between the 1982 problem page and later renumberings remains a separate
source question; the mathematical statement above fixes our quantifiers.

For $\sigma\in\{+1,-1\}$ put
\[
 \delta_\sigma(g)=\norm{g^{-1}\Ric_g-\sigma\Id}_{L^\infty,\op},
 \qquad e_\sigma(M)=\inf_{g\in\Met(M)}\delta_\sigma(g).
\]
Let $P_n^\sigma$ assert the existence of a dimension-uniform positive
threshold for $\delta_\sigma$ implying Einstein existence on the same
smooth manifold. Then
\begin{equation}\label{eq:branches}
                         Q_n\iff P_n^+\lor P_n^-.
\end{equation}
Indeed, $\widehat g=ag$ leaves covariant Ricci unchanged and divides its
endomorphism by $a$. Every nondegenerate interval has a nonzero interior
point, and contains a small interval of its sign. Rescale a sufficiently
small interval about $\sigma$ into that subinterval. Conversely, an
interval about $+1$ or $-1$ supplied by $P_n^\sigma$ proves $Q_n$.

An affirmative answer in all dimensions needs at least one sign in each
dimension. A negative answer needs both signs to fail in one dimension;
the two failures may use different manifolds or sequences of topologies.
A fixed $M$ with $e_\sigma(M)=0$ and no Einstein metric of any sign
disproves that sign branch. Missing higher-dimensional mechanisms are
therefore an obstacle to an all-dimensional affirmative proof, not an
extra requirement for a complete counterexample in one dimension.

\subsection{Rigidity, construction, and global existence}\label{sub:research-roles}
Integral rigidity excludes approximate inputs on certain topologies.
Construction must realize small error in the original, global
$L^\infty$ operator norm on actual smooth closed objects. A global
argument must decide every possible Einstein output on the same object.
These tasks can exchange methods and conclusions. An existence theorem,
a new invariant, a counterexample, or an equivalent reduction is relevant
whenever it returns to \eqref{eq:branches}; the current division into
three subjects imposes no restriction on future routes.

\subsection{Norms, scaling, and the closed category}\label{sub:conventions}
In dimension four write $s=\tr_g\Ric$ and
$E^\circ=\Ric-sg/4$. Weyl norms are Hilbert--Schmidt operator norms on
$\Lambda^2_\pm$; $\Rm_{\rm full}$ has the four-index tensor norm. Thus
\begin{align}
8\pi^2\chi&=\int_M\left(|W^+|^2+|W^-|^2+
                   \frac{s^2}{24}-\frac{|E^\circ|^2}{2}\right)\dd,
                                                        \label{eq:CGB}\\
12\pi^2\tau&=\int_M(|W^+|^2-|W^-|^2)\dd,\label{eq:signature}\\
\int_M|\Rm_{\rm full}|^2\dd&=32\pi^2\chi+
                                      4\int_M|E^\circ|^2\dd.
                                                        \label{eq:full-energy}
\end{align}
The associated Weyl four-tensor has squared norm $4|W^\pm|^2$.
The Eguchi--Hanson energy is $12\pi^2$ in curvature-operator norm and
$48\pi^2$ in full tensor norm. If $\Ric_g\approx-3g$, then
$\widehat g=3g$ satisfies
$\delta_-(\widehat g)=\norm{g^{-1}\Ric_g+3\Id}_\infty/3$.
An integral or weighted smallness statement, or a small-volume error
region, does not imply the required global supremum bound. Complete
noncompact examples do not settle the question in the closed category.

\section{Relative integral Einstein defect}\label{sec:energy}
\subsection{From exact gaps to integral gaps}\label{sub:analytic-change}
The preserved baseline established an exact Einstein gap under a
metricwise monopole projection lower bound and a relative Euler budget.
Its proof used regularity for exact Einstein geometry, fixed-scale local
developments allowing collapse, and an excluded curvature spectrum.
Approximate versions had additional local regularity or degeneration
hypotheses. The first round examined local spectral polynomials and
invariant scalar multipliers. Its sign obstruction and Cotton examples
showed why small Ricci error could not simply be differentiated or fed
into a scalar maximum principle. Round 2 then built a global spinor Weyl
model and established first-order orbit coercivity and weak Bianchi
control. The remaining obstacle was an endpoint product: a bounded
$W^{1,2}$ field with an $L^2$-small remainder did not provide enough
regularity to test the divergence against itself.

Rounds 3--4 changed the analytic mechanism: an actual bounded
spinor on the original metric supplies the derivative control needed for
an integral contradiction. First the covariant Hessian estimate closed
the argument under pointwise Ricci control. Then scalar clipping
preserved the spinor amplitude and the topology-sensitive moment deficit
using only quadratic Ricci error. This removed the remaining pointwise
hypothesis. Round 5 already extended that mechanism to
the polarization capacity of actual special monopole classes.

This excludes candidate approximate inputs, rather than constructing
Einstein metrics. To prove a branch of \eqref{eq:branches}, one would
still need a sufficiently exhaustive topological exclusion or an
independent existence theorem for the remaining cases. The corresponding
proofs are \roundref{round3_rigidity.md}{round 3},
\roundref{round4_rigidity.md}{round 4}, and
\roundref{round5_rigidity.md}{round 5}.

\subsection{Exact optimization over the Einstein constant}\label{sub:energy-optimization}
Set $V=\Vol_g(M)$, $R_2=\int_M|\Ric|^2\dd$, $S=\int_Ms\dd$, and
\begin{equation}\label{eq:energy-definition}
 \cE(g)=\inf_{\lambda\in\R\setminus\{0\}}
       \frac{\int_M|\Ric-\lambda g|^2\dd}{\lambda^2V}.
\end{equation}
\begin{proposition}\label{prop:optimization}
If $R_2>0$, then
\begin{equation}\label{eq:energy-fraction}
 \cE(g)=4-\frac{S^2}{VR_2}
 =\frac{4\int|E^\circ|^2\dd+\int(s-\bar s)^2\dd}{R_2},
 \qquad \bar s=S/V.
\end{equation}
For $S\ne0$ the minimizer is $\lambda=R_2/S$. If $S=0$ and $R_2>0$,
the infimum is $4$ and is not attained at a finite nonzero $\lambda$.
For a Ricci-flat metric the original definition also gives $4$.
\end{proposition}
\begin{proof}
Writing $t=1/\lambda$ transforms the quotient into
$R_2t^2/V-2St/V+4$. Complete the square and use
$R_2=\int|E^\circ|^2+\frac14\int s^2$.
If $R_2=S=0$ the quotient equals $4$ for every allowed $\lambda$.
\end{proof}
Thus zero defect characterizes Einstein metrics with a nonzero Einstein
constant, not all Einstein metrics. The positive monopole-projection
hypotheses below exclude Ricci-flat metrics separately. The functional
differs from the absolute Ricci-energy infimum studied in
\cite{LeBrun01}: a positive lower bound on total Ricci energy does not
give a positive proportion of non-Einstein energy.

\section{The integral calibration theorem and its applications}\label{sec:theorems}
\subsection{The actual-class calibration gap}\label{sub:core-theorem}
For a closed oriented smooth four-manifold put
\begin{gather*}
 E=\Ric+g,\quad e=s+4,\quad
 \mathcal O_0=\Orb\operatorname{diag}(-1/6,1/12,1/12),\\
 F=\dist_{\HS}(W^+,\mathcal O_0)^2,
 \quad\mathcal C=\int_M(|E|^2+F)\dd.
\end{gather*}
The orbit is taken under oriented orthonormal changes of frame; no
global eigenframe is presumed. An actual monopole class is a real first
Chern class of a spin$^c$ structure whose unperturbed Seiberg--Witten
equations have a solution for every smooth metric. For the analytic
argument, it suffices that the associated rescaled equations are solvable
at the metric under study for every smooth positive weight; an actual
monopole class supplies this by conformal covariance. Let $a$ be such a
class and set
\begin{gather}
 p=(a_g^+)^2>0,\qquad q=-s+\sqrt6|W^+|,\qquad q_+=\max(q,0),\notag\\
 \Delta_+=\int_M q_+^2\dd-72\pi^2p,\qquad
 \kappa=\frac{\mathcal C+\Delta_+}{V},\qquad
 B_a=\frac{(-a^2)_+}{V}.\label{eq:defects}
\end{gather}
Here $a_g^+$ is the self-dual harmonic projection and $x_+=\max(x,0)$.

\begin{theorem}[Integral calibration gap, round 4]\label{thm:core}
There is a universal $c>0$ such that
\begin{equation}\label{eq:core-gap}
                         (1+B_a)\kappa\ge c.
\end{equation}
No pointwise curvature or scalar-sign condition, Ricci lower bound,
volume or diameter bound, noncollapse, Sobolev constant, or
curvature-derivative bound is assumed.
\end{theorem}
The actual equation and the class-square factor are hypotheses of the
result. If $w$ is the least Weyl eigenvalue, then
$w\ge-\sqrt{2/3}|W^+|$, and the actual-monopole inequalities are
\begin{equation}\label{eq:spectral-input}
 \int_M(2s/3+2w)_-^2\dd\ge32\pi^2p,
 \qquad \int_Ms_-^2\dd\ge32\pi^2p.
\end{equation}
Since $(-(2s/3+2w))_+\le(2/3)q_+$, they imply $\Delta_+\ge0$.
The negative part is indispensable when scalar curvature changes sign.
The proof of Theorem~\ref{thm:core} occupies
Sections~\ref{sec:hessian}--\ref{sec:self-test}. The rescaled-equation
input is derived in \roundref{round3_monopole_input.md}{the monopole input
note}, building on \cite{LeBrun01,LeBrunConvex}.

\subsection{Topology transfer and the original norm}\label{sub:topology-transfer}
Write $T=2\chi+3\tau$. Suppose each metric admits a selected actual class
with $(a_g^+)^2\ge A>0$ and either
\begin{equation}\label{eq:square-hypotheses}
 a^2\ge-C_AA\qquad\hbox{or}\qquad a^2\ge T-C_AA,
 \qquad C_A\ge0.
\end{equation}
The constant depends only on $C_A$. In particular special classes,
defined by $a^2\ge T$, supply a universal version. The negative-normalized
variables are
\[
 r=\avg{|\Ric+g|^2},\qquad d=(T/A-2/3)_+.
\]
For completeness, the transfer estimates can be written without a
scalar-sign restriction. Set
\begin{gather*}
 \alpha=\sqrt{2/3}|W^+|,\quad w=\lambda_{\min}(W^+),\quad
 t=s+8\sqrt6|W^+|,\quad b=2s/3+2w,\\
 G=(2s/3-2\alpha)^2-b_-^2\ge0.
\end{gather*}
Using \eqref{eq:spectral-input} and
\begin{equation}\label{eq:T-budget}
 4\pi^2T=\int_M\left(\frac{s^2}{24}+2|W^+|^2-
                                      \frac12|E^\circ|^2\right)\dd,
\end{equation}
the identities
\[
 t^2+18G=216(s^2/24+2|W^+|^2)-18b_-^2,
 \qquad27(s^2/24+2|W^+|^2)-q^2=t^2/8
\]
give
\begin{align}
 \int(|E|^2+e^2+t^2+G)\dd
       &\le109\int|E|^2\dd+864\pi^2Ad,\notag\\
 F&\le e^2+t^2+G,\notag\\
 \kappa&\le\frac{245}{2}r+972\pi^2(A/V)d,
 \qquad \frac AV\le\frac{4+r}{4\pi^2}.\label{eq:transfer}
\end{align}
The pointwise estimate for $F$ separates $s\le-3$ and $s>-3$.
For $r\le1$, \eqref{eq:transfer} yields
$\kappa\le1215(r+d)$. Under the first condition in
\eqref{eq:square-hypotheses}, $B_a\le5C_A/(4\pi^2)$. Under the second,
first restrict to $r\le1/4$; then
\[
 4\pi^2T\ge\left(\frac{(4-2\sqrt r)^2}{24}-\frac r2\right)V
                      \ge V/4>0,
\]
so the first condition follows. The complementary $r>1/4$ branch is
already separated from zero. Hence Theorem~\ref{thm:core} gives a
positive bound for $r+d$.

The positive normalization uses classical inequalities alone. With
$r_+=\avg{|\Ric-g|^2}$, \eqref{eq:spectral-input} implies
$A/V\le r_+/(8\pi^2)$. If $r_+\le1/4$, then
$4\pi^2T/V\ge1/4$; if also $d\le1/4$, the preceding upper bound and
$T/A\le11/12$ instead give $4\pi^2T/V\le11/96$.
Thus $r_++d>1/4$. Rescale each sign and apply
Proposition~\ref{prop:optimization} to obtain
\begin{equation}\label{eq:general-gap}
                  \cE(g)+(T/A-2/3)_+\ge c(C_A)>0.
\end{equation}
Details are in \roundref{round4_rigidity.md}{round 4, Sections 7--8}
and \roundref{round4_all_sign_gap.md}{the all-sign corollary}.

\begin{corollary}[Blow-up family]\label{cor:blowups}
Let $X$ be a closed minimal complex surface of general type with
$b_+(X)>1$, $A=c_1^2(X)>0$, and
$M=X\#k\overline{\CP}^{\,2}$. There is $c_*>0$, independent of
$X,k,g$, such that
\begin{equation}\label{eq:blowup-gap}
                      \cE(g)+(1/3-k/A)_+\ge c_*.
\end{equation}
\end{corollary}
\begin{proof}
The full-sign blow-up classes have square $A-k=T$. Averaging their
self-dual projection squares supplies, for each metric, an actual class
with $(a_g^+)^2\ge A$. Apply \eqref{eq:general-gap} with $C_A=0$ in
the second square condition.
\end{proof}
For either sign, evaluating \eqref{eq:energy-definition} at
$\lambda=\sigma$ gives $\cE(g)\le4\delta_\sigma(g)^2$. Consequently
\[
 4\delta_\sigma(g)^2+d\ge c_*,\qquad d=(1/3-k/A)_+.
\]
If $d\le c_*/2$, then $\delta_\sigma(g)\ge\sqrt{c_*/8}$.
An Einstein metric forces $d\ge c_*$, since the scalar monopole
inequality excludes the Ricci-flat exception. Taking width $c_*/2$
gives a closed nonexistence band without confusing strict endpoints.
The constants are not asserted to be optimal. The earlier conditional
smoothing principle, denoted (SC) in the baseline, has been bypassed;
it has not been proved by this argument.

\section{Bounded sections and intrinsic second derivatives}\label{sec:hessian}
\subsection{Covariant integration by parts}\label{sub:hessian-identity}
Let $v$ be a smooth section of a Hermitian bundle with a compatible
connection over a closed four-manifold. Write
$K=\norm v_\infty$, $H=\nabla^2v$, and
$\mathscr F_{ij}=[\nabla_i,\nabla_j]$ on the bundle. Curvature norm means
$|\mathscr F|^2=\sum_{ij}|\mathscr F_{ij}|_{\HS}^2$, with ordered indices.
Integrating the divergence of
$|\nabla v|^2\langle v,\nabla v\rangle^\sharp$ gives
\begin{equation}\label{eq:interpolation}
                       \norm{\nabla v}_4^2\le4K\norm H_2.
\end{equation}
Indeed, the differentiated gradient factor contributes at most
$2K|H||\nabla v|^2$ and the trace Hessian at most
$2K|H||\nabla v|^2$. This uses the amplitude bound, not a uniform
Sobolev constant on a possibly degenerating family.

Commuting covariant derivatives and integrating once yields
\begin{align}
 \int|\nabla^2v|^2\dd={}&\int|\nabla^*\nabla v|^2\dd
   -\int\Ric_{ij}\langle\nabla_iv,\nabla_jv\rangle\dd\notag\\
 &+\sum_{ij}\int\langle\mathscr F_{ij}\nabla_jv,\nabla_iv\rangle\dd
   +\frac12\sum_{ij}\int|\mathscr F_{ij}v|^2\dd.
                                                        \label{eq:hessian-identity}
\end{align}
Real parts of Hermitian products are understood. H\"older and
\eqref{eq:interpolation} bound the two curvature-gradient terms by
$4K(\norm{\Ric}_2+\norm{\mathscr F}_2)\norm H_2$.
Young's inequality therefore gives
\begin{equation}\label{eq:hessian-bound}
 \norm{\nabla^2v}_2^2\le2\norm{\nabla^*\nabla v}_2^2
             +32K^2\norm{\Ric}_2^2+33K^2\norm{\mathscr F}_2^2.
\end{equation}
No derivative of curvature occurs and no curvature smallness is needed.
The complete commutator calculation and independent check are preserved
in \roundref{round3_rigidity.md}{round 3} and
\roundref{round3_rigidity_second_review.md}{its second review}.

\subsection{The positive-spin curvature budget}\label{sub:spin-curvature}
The curvature of the positive spin connection involves $W^+$, scalar
curvature, the relevant trace-free Ricci block, and the determinant
connection. The $W^-$ spin factor acts trivially. This representation
fact does not declare $W^-$ irrelevant to the whole geometry; it identifies
the curvature actually needed in \eqref{eq:hessian-bound}.

The rescaled monopole equation controls $F_A^+$. Chern--Weil supplies
the other half through the real first Chern class:
\begin{equation}\label{eq:determinant-budget}
 \norm{F_A}_2^2=2\norm{F_A^+}_2^2-4\pi^2a^2
             \le4\int f^2\dd+4\pi^2(-a^2)_+,
\end{equation}
where the last inequality uses the spinor bound obtained below.
This explains the negative-square cost. It is a consequence of the
actual determinant curvature, not a removable notation choice.

\section{A scalar-safe weight and near-equality information}\label{sec:weight}
\subsection{Truncation with a quadratic cost}\label{sub:clipping}
For every smooth positive $f$, solve
\begin{equation}\label{eq:rescaled-SW}
 D_A\Phi=0,\qquad -iF_A^+=f\sigma(\Phi),\qquad
 \psi=2\sqrt2\,\sigma(\Phi),\quad
 \rho=|\Phi|^2=|\psi|,\quad u=3\rho/2,
\end{equation}
with $|\sigma(\Phi)|^2=|\Phi|^4/8$. One obtains existence by solving the
ordinary equations at $\widehat g=f^{-2}g$ and using the conformal spinor
weight $\Phi=f^{-3/2}\widehat\Phi$; see \cite{LeBrunConvex}, Lemma 3.7.
Existence here requires an actual monopole class.

Set
\begin{equation}\label{eq:weight}
 Q=\max(q,3),\qquad f_0=\max\{\sqrt Q,(-s)_+/4\}.
\end{equation}
On each fixed compact $(M,g)$ choose smooth majorants $f_j\ge f_0$
converging uniformly to $f_0$, and solve \eqref{eq:rescaled-SW} with
$f=f_j$. The maximum principle gives $f_j\rho\le-s$ at a positive
maximum of $\rho$, hence
\begin{equation}\label{eq:amplitude-bound}
                         0\le\rho\le4,\qquad0\le u\le6.
\end{equation}
The weight has controlled costs
\begin{equation}\label{eq:clipping-cost}
 0\le Q^2-q_+^2\le9(s+4)^2,
 \qquad0\le f_0^2-Q\le\frac9{16}(s+4)^2.
\end{equation}
The additional amplitude control therefore consumes only the allowed
scalar $L^2$ error. Direct use of $\int f_j^4$ would ask for an
unavailable fourth moment. Smooth approximation is completed for fixed
$(g,a)$ before taking a diagonal choice along a contradiction sequence;
derivatives of $f_j$ need no uniform bound across that sequence.

\subsection{Actual and ideal moments}\label{sub:moments}
Subtracting the spinor and self-dual-form Weitzenb\"ock formulas produces
four nonnegative defects:
\begin{align}
 D_{\rm amp}&=2\int|d\rho|^2\dd,\notag\\
 D_{\rm quad}&=\int(4\rho|\nabla_A\Phi|^2-|\nabla\psi|^2)\dd,\notag\\
 D_H&=\int(|d\psi|^2+|d^*\psi|^2)\dd,\notag\\
 D_{\rm eig}&=2\int\left(\langle W^+\psi,\psi\rangle+
                                 \sqrt{2/3}|W^+|\rho^2\right)\dd.
                                                        \label{eq:four-defects}
\end{align}
The second follows from the quadratic spinor-map derivative estimate,
the fourth from the least-eigenvalue bound. Their sum $D_0$ satisfies
\begin{equation}\label{eq:exact-defect}
                   \frac{27}{8}D_0=\int(qu^2-f_ju^3)\dd.
\end{equation}
Define the ideal moments
$C=\int Qu^2\dd$, $Z=\int\sqrt Q\,u^3\dd$, and
$E_Q=\int Q^2\dd$. Then
\[
 P:=C-Z=\frac{27}{8}D_0+
       \int(Q-q)u^2\dd+\int(f_j-\sqrt Q)u^3\dd\ge\frac{27}{8}D_0.
\]
Topology controls a different quantity:
\begin{equation}\label{eq:actual-moment}
 C_{\rm act}=\int f_j^2u^2\dd=18\norm{F_A^+}_2^2\ge72\pi^2p.
\end{equation}
The comparison needed to transfer that bound is
\[
 0\le C_{\rm act}-C\le\frac{81}{4}\int(s+4)^2\dd+o(1).
\]
By H\"older, $C\le E_Q^{1/3}Z^{2/3}$; since $Z\le C$,
\[
 0\le E_Q-C\le S_0+o(1),\qquad
 S_0=\Delta_++\frac{117}{4}\int(s+4)^2\dd.
\]
With $\theta=(Z/E_Q)^{1/3}$ and $H_0=E_Q^{1/3}Z^{2/3}-C$,
\[
 E_Q-C=H_0+\frac{1+\theta}{\theta^2}(P+H_0)\ge2P.
\]
The positive projection excludes the identically zero spinor, so these
moment expressions are legitimate. Applying the corresponding moment
identity to $Q^2\dd/E_Q$ gives
\begin{align}
 D_0&\le\frac4{27}S_0+o(1),\notag\\
 \int Q(u-\sqrt Q)^2\dd&\le S_0+o(1),\notag\\
 \norm{d\rho}_2^2,\ \norm{d\psi}_2^2&\le\frac2{27}S_0+o(1).
                                                        \label{eq:moment-estimates}
\end{align}
These control amplitude, derivatives, and the eigenline simultaneously.
Every $o(1)$ is first made arbitrarily small on the fixed metric, and
small relative to $V$ when metrics vary. No convergence of solutions as
$j\to\infty$ is asserted. The complete algebra is in
\roundref{round4_rigidity.md}{round 4, Sections 3--4}.

\section{The smooth Weyl model and the Bianchi self-test}\label{sec:self-test}
\subsection{A polynomial model through the spinor zero set}\label{sub:Weyl-model}
Use the global smooth tensor
\begin{equation}\label{eq:model}
                 A_\psi=\frac{\rho^2}{24}\Id-\frac18\psi\otimes\psi.
\end{equation}
At $\rho^2=2$ its spectrum is $(-1/6,1/12,1/12)$; at $\rho=0$
the same polynomial remains smooth. There is no division by amplitude.

The error ledger gives $\int(q-9/2)^2\dd\le12\mathcal C$;
clipping at $3$ does not increase this distance. Thus
\eqref{eq:moment-estimates} controls $\rho-\sqrt2$ in $L^2$.
Choose a measurable least-eigenline projection $\Pi$. On
$\{F\le1/256\}$ its spectral gap is at least $1/8$, while the
complement has volume at most $256\int F\dd$. The eigenline defect
and $\rho\le4$ control $(I-\Pi)\psi$ globally. For
$A_0=I/12-\Pi/4$,
\[
 |A_\psi-A_0|^2\le\frac{\rho^2}{16}|(I-\Pi)\psi|^2+
                                      \frac{(\rho^2-2)^2}{48}.
\]
Consequently
\begin{equation}\label{eq:model-error}
 R=W^+-A_\psi,\qquad
                   \norm R_2^2\le2\Delta_++8500\mathcal C+o(1).
\end{equation}
The coefficient $8500$ is sufficient, not optimized.

Apply \eqref{eq:hessian-bound} to $\Phi$. Its rough Laplacian is
controlled by $s$, $f_j$, and \eqref{eq:amplitude-bound}; its connection
curvature is controlled by \eqref{eq:determinant-budget} and the positive
spin representation. Hence
\begin{equation}\label{eq:second-derivatives}
 \int\left(|\nabla_A^2\Phi|^2+|\nabla_A\Phi|^4+
                                  |\nabla^2A_\psi|^2\right)\dd
       \le C\bigl(V+\mathcal C+(-a^2)_+\bigr)+o(1).
\end{equation}
The last term uses the quartic spinor polynomial in \eqref{eq:model}.
All terms are unnormalized before division by $V$.

\subsection{Incompatible divergence-energy bounds}\label{sub:contradiction}
Let $\mathcal W$ denote the identification with a Weyl four-tensor, and
$J_\psi=\diver\mathcal W(A_\psi)$. Contracted Bianchi,
$W^+=A_\psi+R$, and integration by parts give, for smooth test tensors $U$,
\begin{equation}\label{eq:weak-Bianchi}
 \left|\int\langle J_\psi,U\rangle\dd\right|
       \le C(\norm E_2+\norm R_2)\norm{\nabla U}_2.
\end{equation}
This is an intrinsic negative-order estimate for the current metric.
Equation~\eqref{eq:second-derivatives} permits $U=J_\psi$, and therefore
\begin{equation}\label{eq:J-upper}
 \avg{|J_\psi|^2}\le C\sqrt\kappa\sqrt{1+\kappa+B_a}+o(1).
\end{equation}
The weighted Hodge identity, including amplitude derivatives, is
\begin{align}
 \int|J_\psi|^2\dd={}&-\frac1{64}\int\rho^2B_\psi\dd
          +\frac1{16}\int\rho^2|d\psi|^2\dd
          -\frac5{96}\int\rho^2|d\rho|^2\dd\notag\\
 &+\frac18\int\rho\langle\iota_{\nabla\rho}\psi,*d\psi\rangle\dd,
                                                        \label{eq:weighted-Hodge}\\
 B_\psi={}&(s/3)\rho^2-2\langle W^+\psi,\psi\rangle.\notag
\end{align}
The curvature contribution equals
\[
 -\frac{\rho^2B_\psi}{64}
 =-\frac{s\rho^4}{192}-\frac{\rho^6}{384}
                    +\frac{\rho^2\langle R\psi,\psi\rangle}{32}.
\]
At $s=-4$, $\rho^2=2$, and $R=0$ it is $1/16$.
The four defects control the differential and mixed terms, so
\begin{equation}\label{eq:J-lower}
 \avg{|J_\psi|^2}\ge\frac1{16}-C(\sqrt\kappa+\kappa)+o(1).
\end{equation}
If $(1+B_a)\kappa\to0$, the upper bound tends to zero while the lower
bound tends to $1/16$. This proves Theorem~\ref{thm:core}.
The first-jet contraction and weighted Hodge calculation are in
\roundref{round3_monopole_input.md}{the monopole-input appendices};
the integral closure is in \roundref{round4_rigidity.md}{round 4,
Sections 5--7}.

A bounded $W^{1,2}$ model, an $L^2$-small remainder, and a small
$H^{-1}$ divergence alone do not control derivative energy. The weighted
identity by itself merely retains the positive energy at calibration.
The additional derivative in \eqref{eq:second-derivatives} is what makes
the Bianchi self-test a contradiction without extracting a regular
limit manifold.

\subsection{A portable conditional criterion}\label{sub:portable-criterion}
\begin{proposition}[Positive-energy self-test]\label{prop:self-test}
Suppose $\kappa,B\ge0$ and a legitimately testable field $J$ satisfy
\[
 \avg{|J|^2}\le C_1\sqrt\kappa\sqrt{1+\kappa+B},\qquad
 \avg{|J|^2}\ge m-C_2(\sqrt\kappa+\kappa),\quad m>0.
\]
Then $(1+B)\kappa\ge c>0$, where $c$ depends only on $m,C_1,C_2$.
\end{proposition}
\begin{proof}
Choose $c\le1$ so that
$C_1\sqrt{c+c^2}\le m/4$ and $C_2(\sqrt c+c)\le m/2$.
If $(1+B)\kappa<c$, then $\kappa<c$ and
$\kappa(1+\kappa+B)<c+c^2$. The upper bound is at most $m/4$
and the lower bound at least $m/2$, a contradiction.
\end{proof}
The implication is elementary. Extending it to another geometric
equation requires an actual solution, a field $J$, an intrinsically
controlled second derivative, and a strictly positive retained energy
$m$. The proposition alone supplies none of those geometric inputs.

\section{Actual class sets and topological scope}\label{sec:classes}
\subsection{Special and ordinary monopoles}\label{sub:special}
A special class satisfies $a^2\ge T$. Nonzero chamber-independent
Seiberg--Witten data and suitable nonzero Bauer--Furuta data supply
generic monopoles under their precise hypotheses
\cite{Kotschick,Ishida,IL}. Generic perturbations have expected dimension
$(a^2-T)/4$, so persistence and transversality imply the special condition.
Metricwise existence for the unperturbed equations alone does not
justify that argument. Ordinary and special sets must remain distinct.

For a finite nonempty set $\mathscr C_0$ of actual classes define
\begin{equation}\label{eq:capacities}
 \beta_0^2=\max_{v\in\Hull(\mathscr C_0)}v^2,\qquad
 \alpha_0^2=\inf_{H\in\Gr^+_{b_+}}
                       \max_{a\in\mathscr C_0}(a_H^+)^2.
\end{equation}
All maximal positive subspaces occur in the infimum, without assuming
that every one is realized by a metric. If $v=\sum_i t_i a_i$ is a
convex combination, then
\[
 v^2\le(v_H^+)^2\le\sum_i t_i(a_{i,H}^+)^2
                                     \le\max_i(a_{i,H}^+)^2.
\]
Thus $\beta_0^2\le\alpha_0^2$. The finite maximum selects a class
for each metric; the Grassmannian infimum need not be attained.

\begin{theorem}[Special-monopole gap, round 5]\label{thm:special}
Suppose $b_+(M)\ge2$ and the set of actual special monopole classes has
$\alpha_{\rm sp}^2>0$. There is a universal $\varepsilon_0>0$ such that
\begin{equation}\label{eq:special-gap}
 \cE(g)+\left(T/\alpha_{\rm sp}^2-2/3\right)_+\ge\varepsilon_0
\end{equation}
for every metric. The same conclusion holds with any positive
$A\le\alpha_{\rm sp}^2$, including $A=\beta_{\rm sp}^2$ when positive.
\end{theorem}
\begin{proof}
For each metric select a class attaining the largest self-dual projection.
It has $p\ge A$ and $a^2\ge T$. Apply \eqref{eq:general-gap}.
Finiteness of the real monopole set for $b_+\ge2$ is the classical
result in \cite{LeBrunConvex}; no bound on its cardinality is needed.
\end{proof}
For all ordinary classes on a fixed manifold one instead keeps
$C_M=\max_a(-a^2)_+/\alpha_{\rm all}^2$ and a constant depending on
$C_M$. Convex geometry alone cannot eliminate it: in signature $(2,1)$,
the vectors $\pm(1,0,L),\pm(1,0,-L)$ have
$\alpha^2=\beta^2=1$ and squares $1-L^2$. They are an algebraic
countermodel to the inference, not a realized smooth monopole family.

\subsection{Connected sums and actual gluing input}\label{sub:connected-sums}
Suppose $X$ supports a relevant nonzero chamber-independent SW class $c$,
$b_+(X)\ge2$, and $A=c^2>0$. Let $b_+(N)=0$ and assume an applicable
actual-monopole gluing theorem supplies the full-sign classes
$a_\epsilon=c+\sum_i\epsilon_i e_i$, with
$e_i e_j=-\delta_{ij}$. Put $k=b_2(N)$, $\ell=b_1(N)$, and
$D=k+4\ell$. Then
\[
 T(X\#N)=T(X)-D,\qquad
 a_\epsilon^2=A-k\ge T(X)-D,
\]
because $A\ge T(X)$. Averaging projected squares supplies $p\ge A$;
thus
\begin{equation}\label{eq:surgery-gap}
 \cE(g)+\left((T(X)-D)/A-2/3\right)_+\ge\varepsilon_0.
\end{equation}
Ishida's Proposition 6 in \cite{Ishida} gives the relevant source input;
the standard family is
$N=k\overline{\CP}^{\,2}\#\ell(S^1\times S^3)$.
See \roundref{round5_rigidity.md}{round 5, Section 4} for the full
hypotheses. Betti numbers and a formal surgery formula do not manufacture
solutions. Graph-block self-gluing is a different operation from a
spherical connected sum; a local circle or sphere does not identify the
global surgery.

\section{Compatibility with actual approximate constructions}\label{sec:compatibility}
\subsection{Ratio one and the critical ratio}\label{sub:ratios}
The slow-twist surface bundles $E_D$ and the specified closed sheared
objects of rounds 11--12 have $\tau=0$ and $\alpha^2=2\chi$.
For the indicated $E_D$ examples, both orientations satisfy
\[
 \alpha_{\rm all}^2=\beta_{\rm all}^2
 =\alpha_{\rm sp}^2=\beta_{\rm sp}^2=T.
\]
Their ratio is $1$, whereas \eqref{eq:special-gap} concerns the region
near $2/3$. Their $e_-=0$ is compatible with the obstruction.
The class computation for the one-block self-gluings is in
\roundref{round12_global.md}{round 12}; their variational infima are
computed in \roundref{round11_variational.md}{round 11}. Attainment of
the full-curvature infimum there is equivalent to the unresolved
Einstein-existence question. It cannot be inferred from an infimum value.

For a fixed resolution of a real-hyperbolic $A_1$ orbifold, assume an
actual uniformly near-negative-Einstein sequence with a fixed
$\Ric=-g$ exterior, limiting volume $V_0$, and replacement-region
volume tending to zero. Set $B=V_0/(6\pi^2)$ and let $m$ count the
positive-signature contributions of the Eguchi--Hanson replacements.
The accounting in \roundref{round5_rigidity.md}{round 5, Section 6} gives
\[
 T=B+6m,\quad\lim\int|W^+|^2\dd=12\pi^2m,
 \quad\alpha_{\rm all}^2\le\frac43B+m,
 \quad T/\alpha_{\rm all}^2\ge\frac34
\]
when the capacity is positive. The $L^1$ Weyl contribution of the
vanishing-volume regions tends to zero.

The filling-independent version requires a fixed smooth $M$, uniform
Ricci defect tending to zero, $V\to V_0$, replacement volume tending
to zero, a fixed negative Einstein orbifold exterior $O$, and convergence
of the exterior first and second moments of $|W^+|$. Gromov--Hausdorff
convergence alone supplies none of the moment assumptions. Write
$w_0=|W_O^+|$ and let $E_{\rm bub}\ge0$ be limiting self-dual Weyl
energy in the replacement regions. Define
\begin{gather*}
 T_0=\frac1{4\pi^2}\int_O(2/3+2w_0^2),\qquad
 J_0=\frac1{72\pi^2}\int_O(4+\sqrt6w_0)^2,\\
 D_{\rm orb}=J_0+\frac{E_{\rm bub}}{12\pi^2}.
\end{gather*}
Then
\begin{gather}
 T=T_0+\frac{E_{\rm bub}}{2\pi^2},\qquad
 \alpha_{\rm all}^2\le\min\{V_0/(2\pi^2),D_{\rm orb}\},\notag\\
 T-\frac23D_{\rm orb}
 =\frac4{9\pi^2}\left[
       \int_O\left(w_0-\frac1{2\sqrt6}\right)^2+E_{\rm bub}\right].
                                                        \label{eq:concentration}
\end{gather}
The bubble contribution has ratio six, so adding concentration does not
move these data toward the critical ratio. Near-criticality requires
vanishing normalized bubble energy, the calibrated bulk Weyl modulus,
and saturation by actual monopole classes simultaneously.

Real-hyperbolic, complex-hyperbolic, and bidisk bulks remain at ratio at
least $3/4$ in both orientations under these hypotheses. A K\"ahler bulk
in its complex orientation has the stronger bound
$1+E_{\rm bub}/V_0$. In the opposite orientation the formal critical
Chern slope is $c_1^2/c_2=32/11$. It is a selection condition, not a
constructed or monopole-saturated example. A constant Weyl norm alone
does not impose the forbidden eigenvalue spectrum. These calculations
are in \roundref{round7_endpoint.md}{round 7's endpoint note}.

\subsection{Joint design on the same smooth object}\label{sub:joint-design}
A construction and obstruction must use the same topology, actual class
set, gluing marking, capacity, and full Ricci norm. If surgery moves
$T/\alpha^2$ toward $2/3$, its Ricci cost and the persistence of the
actual classes must be recomputed together. For example, when
$A=T(X)>0$, \eqref{eq:surgery-gap} prevents connected sums with $D/A$
near or above $1/3$ from retaining arbitrarily tight near-Einstein
metrics. This constrains the modified object, not the original $E_D$.
The first incompatible identity may yield a useful surgery-cost or
class-failure theorem.

\section{Independent uses of the earlier rigidity tools}\label{sec:legacy}
\subsection{Exact gaps, local developments, and failed shortcuts}\label{sub:legacy-local}
The baseline hypothesis $\mathrm{MP}(A)$ means a metricwise actual
monopole projection lower bound $A$; together with $\chi\le LA$, its
exact Einstein gap has a scope not identical to the new selected-class
square bounds. The preserved
\baseref{20260916_rigidity.tex}{rigidity baseline} and
\baseref{20260916_uniform_gap.tex}{uniform-gap baseline} retain
low-energy regions, path lifting, covering multiplicity, and transfer
of averaged defects. The relevant exact-geometry regularity sources
include \cite{CT,CN,JN}. Earlier necessary collapse or diameter escape
statements become vacuous inside a band now excluded outright, but may
still have content under their broader or different hypotheses.

Several analytic distinctions also survive. In the negative double-root
ray $u\operatorname{diag}(-2,1,1)/\sqrt6$, the proposed global invariant
scalar multiplier would require
$(u^{2/3}\varphi'(u))'\ge0$, hence $\varphi'\ge0$, while the required
reaction has the opposite sign. This excludes those sign requirements
for that scalar strategy, not every maximum principle. A local spectral
polynomial still excludes a sufficiently small spectral box when
$s=-4$ and $\diver W^+=0$.

An actual oscillatory metric perturbation of size $\epsilon/k^2$ on a
closed negative Einstein background has Ricci error $O(\epsilon)$
and bounded full curvature but Cotton energy at least
$c\epsilon^2k^2$. Thus small Ricci error need not control Cotton energy.
Separate algebraic examples show that quadratic calibration defect
cannot automatically control cubic Weyl concentration; they are not
claimed to be metric curvature examples. Bounded fields and small
$L^2$ remainders likewise do not control the endpoint product
$R\nabla a$ in $H^{-1}$. The actual spinor Hessian estimate supplies
the missing information. See \roundref{rigidity.md}{the first rigidity
note} and \roundref{round2_rigidity.md}{round 2}.

For each fixed full-curvature bound $K$, the independent geometric result
\[
 |\Rm|\le K\quad\Longrightarrow\quad
 \avg{|\Ric+g|^2+\dist(W^+,\mathcal O_0)^2}\ge c(K)>0
\]
allows collapse. Its proof uses buffered exponential developments,
multiplicity bounds, and elliptic estimates. On already controlled
charts the Weyl-divergence symbol satisfies
$|\sigma_D(\xi)A|^2=|\xi|^2|A|_{\HS}^2$ and permits strong local
$L^2$ transfer without an $L^2$ Cotton bound; see
\roundref{round2_weak_curvature.md}{the weak-curvature note}.
The localized degeneration conclusion concerns fixed-radius
neighborhoods of high-curvature sets occupying asymptotically all volume,
not positive volume of the high-curvature set itself.

For comparison, the exact Ricci-box estimate for eigenvalues
$-1+a_i$, $|a_i|\le d\le1$, is
\[
 \frac{s^2}{24}-\frac12|E^\circ|^2\ge
 \begin{cases}
 \frac23(1-d)^2,&0\le d\le1/3,\\
 \frac23(1+d)(1-2d),&1/3\le d\le1.
 \end{cases}
\]
Thus $d<1/2$ implies strict Hitchin--Thorpe. Connector-volume bounds
require a genuine match between the retained end volume and the closed
Chern--Weil budget. For the fixed blow-up ratio $p=k/A>1/3$ the
algebraic quotient optimization yields
\[
 d\ge\begin{cases}
 \sqrt{(3p-1)/(p+5)},&1/3<p\le2\sqrt3-3,\\
 [\sqrt{(3p-1)(3p+15)}-(3p-1)]/8,&2\sqrt3-3\le p<1,
 \end{cases}
\]
and $d\ge1/2$ for $p\ge1$. These are algebraic optima, not claims of
metric attainment; they vanish at the critical ratio and do not replace
\eqref{eq:blowup-gap}. See \roundref{curvature_box.md}{the curvature box}
and \roundref{fixed_ratio.md}{the fixed-ratio calculation}.

\subsection{Localized flow as an independent possible tool}\label{sub:legacy-flow}
The baseline flow calculations retain support buffers, the choice
$a=\chi^2$ with $|da|^2/a=4|d\chi|^2$, signed Hessian integration by
parts, shifted curvature $\mathsf T=\Rm+G$ and $U=a\mathsf T$,
energies at all exponents, critical-to-subcritical Moser estimates,
first-exit estimates on common supports, and terminal calibration.
Highest cutoff derivatives, background-connection control, a common
regularization time, and existence of the limiting flow require separate
proofs. This remains a possible tool for objects beyond the integral
argument. It is not a prerequisite for the completed gap theorem.

\section{Extending the topology covered by the integral theorem}\label{sec:cost-directions}
\subsection{Economical metric-dependent class selection}\label{sub:class-selection}
\begin{problem}\label{prob:selection}
For a finite set $\mathscr C$ of actual classes and a target $A>0$, define
\begin{equation}\label{eq:Theta}
 \Theta_A(g)=\min\left\{\frac{(-a^2)_+}{A}:a\in\mathscr C,
                                      (a_g^+)^2\ge A\right\}.
\end{equation}
When the eligible set is nonempty, can actual class geometry bound
$\Theta_A(g)$ uniformly, or sufficiently along near-critical sequences?
\end{problem}
This asks less than bounding the most negative square of every class.
Finite surgery families, signed Bauer--Furuta families, and symplectic
classes in both orientations offer finite projection-selection problems.
One must first verify that the proposed classes are actual. A bound
depending on $M$ yields a per-manifold result; a family-uniform conclusion
needs control as the topology changes. Exact linear algebra can solve
the finite selection step but cannot prove metric realization of every
formal positive subspace.

\subsection{Can another positive energy absorb negative-square growth?}\label{sub:cost-compensation}
The cost $B_a$ measures determinant anti-self-dual energy. One may ask
whether large $B_a$ forces an additional topological or spatial
phase-separation cost for actual solutions. A possible target is an
inequality $\kappa\Xi(B_a)\ge c$ for which near-critical data control
$\Xi(B_a)$, or a second contradiction when $B_a\to\infty$.
There is no reason for $F_A^-$ to become small automatically. A candidate
joint inequality must survive both the formal convex countermodel in
Section~\ref{sub:special} and tests on actual known monopole families.

\subsection{Effective widths and geometric sharpness}\label{sub:effective-constants}
Tracing every sufficient constant in \eqref{eq:J-upper} and
\eqref{eq:J-lower} can produce a computable, possibly tiny, gap useful
for finite examples and reproduction. The more structural question is
whether its smallness reflects estimation loss or a sequence of actual
objects approaching equality. Optimizing permitted spectra, amplitudes,
Hodge energy, and topology jointly may address this. An optimal matrix
calculation alone is not a sharp geometric constant.

\subsection{Boundary and localized relative obstructions}\label{sub:boundary}
Closedness removes the boundary terms in the self-test. A connector
problem instead needs a domain or cutoff version with explicit
$\nabla\chi$, $\nabla^2\chi$, and actual boundary Cauchy data.
A concrete objective is a relative self-test identity with a computed
flux, followed by a determination of whether surrounding global geometry
controls that flux. A total energy bound does not ensure a region with
both small cutoff errors and the correct topological information.
No boundary-free application of \eqref{eq:core-gap} to a connector is
justified by the present theorem.

\section{Other geometric mechanisms and coupled equations}\label{sec:other-mechanisms}
\subsection{Realizable calibration orbits}\label{sub:other-orbits}
Replacing $\mathcal O_0$ requires three new inputs: an actual equation
or topological inequality forcing the candidate orbit, a global smooth
replacement tensor, and a positive Hodge--Bianchi self-test energy.
Negative K\"ahler--Einstein products, real-hyperbolic metrics, and
complex-hyperbolic metrics are essential controls. An orbit already
realized by an Einstein metric cannot produce the same universal
contradiction.

The product-like geometry of the ratio-one sheared examples should
instead motivate a global transition question. Can monodromy or shear
force a nonzero flux or cost when locally permitted geometries are
assembled? Such a result could link integral estimates with global
topology without forcing the examples into the $2/3$ calibration.

\subsection{Several actual classes and a common differential system}\label{sub:multiple-classes}
For monopole-induced forms $\psi_\alpha$, consider their Gram matrix,
projection fields, and common covariant derivatives. Simply summing
scalar Weitzenb\"ock inequalities can cancel the dependence on $b_+$.
New information would have to come from simultaneously realizable
phases, cross derivatives, or transfer through zero sets.

A testable aim is to derive the equation for a joint projection tensor
$P$, keeping all sources in $\nabla P$, and ask whether topology-forced
rank changes have a quantitative cost. Products, K\"ahler--Einstein
metrics, and the actual sheared families test whether that cost can
concentrate and disappear. No such joint-rank obstruction is proved here.

\subsection{Source defects and nonabelian systems}\label{sub:sources-nonabelian}
For perturbed Dirac, curvature, or Hodge equations, every source must
enter its correct dual norm. An $H^{-1}$ source needs a controlled
$H^1$ test field; an $L^2$ source cannot be differentiated without a
further bound. An explicit source version of
Proposition~\ref{prop:self-test}, followed by an actual geometric system
supplying its budgets, is a concrete extension problem.

The nonabelian work has partial results that identify what remains.
\roundref{round8_nonabelian.md}{Round 8} gives a Chern--Weil lower bound
for projective-line curvature or second-fundamental-form energy.
\roundref{round9_nonabelian.md}{Round 9} constructs anisotropic seeds
solving the Dirac and determinant-trace equations, leaving trace-free
curvature unsolved. \roundref{round10_nonabelian.md}{Round 10} shows a
small-residual obstruction for balanced full-rank seeds with small
anti-self-dual curvature. \roundref{round11_rankone.md}{Round 11} retains
the line second-fundamental-form term at rank one, without forcing the
Dirac right inverse to become singular.

\roundref{round12_uniformization.md}{Round 12} corrects the determinant
central term in an explicit Clifford-block connection and tracks graded
energy. An indefinite Hermitian connection cannot use a positive-unitary
energy argument without that extra identity. The resulting classification
is conditional on solving the explicitly defined system. Ordinary SW
classes are not known to supply those coupled solutions, and the system
has not been identified with every convention in published nonabelian
uniformization converses. The earlier literal rank-one calculation in
\roundref{round7_rigidity.md}{round 7} has the same convention limitation.
The \roundref{round12_control.md}{Schoen-surface control} shows that
negative Einstein geometry, ordinary SW information, and signature zero
do not alone force local symmetry. Moduli invariants or deformation
paths supplying the missing existence input remain possible directions.

\section{Positive Ricci pinching and higher dimensions}\label{sec:positive-higher}
\subsection{The positive branch}\label{sub:positive}
Tight positive Ricci pinching gives positive scalar curvature. The usual
nonzero monopole projection therefore is not an available input for
these candidate manifolds. One can investigate harmonic two-forms of a
putative Einstein output, Gram tensors, curvature actions, or spectral
thresholds, but the identities must be compatible with actual Einstein
metrics on $\CP^2$, suitable del Pezzo surfaces, and $S^2\times S^2$.
It must first be shown that an identity retains topological rank
information before it can yield a rank obstruction.

When $b_+=1$, chamber choice, metric dependence of classes, and wall
crossing require a new argument. The all-positive-weight existence used
above cannot be transferred from the high-$b_+$ setting without proof.
A theorem in a specified chamber may have value while leaving the full
metric quantifier in $P_n^+$ untouched.

\subsection{Actual systems in higher dimensions}\label{sub:higher}
For $n$ dimensions and $p>0$ define
\begin{equation}\label{eq:higher-energy}
 \cE_{n,p}(g)=\inf_{\lambda\ne0}
 \frac{\int_M|\Ric-\lambda g|^p\,d\mu_g}{|\lambda|^p\Vol_g(M)}.
\end{equation}
This is invariant under constant scaling with the corresponding change
of $\lambda$. The exponent $p=n/2$ is curvature-critical, but choosing it
does not provide a higher-dimensional substitute for actual monopoles.
Possible routes involve Dirac indices, special structures, codimension-four
ALE families, or genuinely solvable coupled gauge equations. Each route
needs existence, bounded amplitude, second-derivative estimates, and a
positive test quantity. The exponent matching in the four-dimensional
$L^2$ Hessian absorption must be rederived.

An independent affirmative route could prove Einstein existence for all
remaining candidates or on covers with a valid descent mechanism.
It need not encode every topology in an obstruction theorem. These
existence and nonexistence routes remain open alternatives within the
same objective.

\section{Three concrete mathematical experiments}\label{sec:experiments}
\subsection{Comparing the costs of actual class choices}\label{sub:experiment-A}
Writing $D=\mathcal C+\Delta_+$, the core theorem is
\[
          \left(1+\frac{(-a^2)_+}{V}\right)\frac DV\ge c.
\]
In the fixed negative normalization, follow the four ratios
$D/V$, $(-a^2)_+/V$, $A/V$, and $T/A$ separately along a proposed
counterexample. If the first tends to zero while the second diverges,
investigate whether a different actual class avoids the divergence,
whether determinant energy must concentrate, or whether another bounded
test section is needed.

For an explicit finite signed family and a fixed positive subspace $H$,
compute every $(a_H^+)^2$ and $a^2$, then select high projection at low
negative-square cost. Only afterwards seek uniform control as $H$
ranges over actual metric polarizations. A formal sequence $H_i$ with
divergent selection cost is not yet a geometric counterexample: metric
realization is a separate problem. Either a broader uniform gap or an
actual failure of capacity-only stability would resolve a new issue.

\subsection{A local spectral criterion retaining all sources}\label{sub:experiment-B}
The early scalar multiplier, Cotton, cubic concentration, and endpoint
product examples in Section~\ref{sub:legacy-local} show what cannot be
discarded. For a new actual equation, first derive the first-order symbol
of its orbit restriction and then its full second-order error equation.
With the convention $R=W^+-A$, the schematic identity must read
\begin{equation}\label{eq:source-identity}
 \diver\mathcal W(A)
       =-\diver\mathcal W(R)+\mathcal T(E)+\mathcal S.
\end{equation}
Here $\mathcal T$ is the explicitly derived first-order Ricci operator;
$\mathcal S$ must come from the actual auxiliary equation. The minus
sign follows from $W^+=A+R$. Any different sign requires redefining $R$.
Writing $J=\diver\mathcal W(A)$, a useful sufficient target is
\begin{equation}\label{eq:source-absorption}
 |\langle\mathcal S,J\rangle|\le o(V)+\theta\norm J_2^2,
                      \qquad 0\le\theta<1,
\end{equation}
where the pairing is the appropriate integral or dual pairing. This is
a proposed estimate, not a consequence of an unspecified weak norm.
If the source-free Bianchi estimate and \eqref{eq:source-absorption}
both hold, the factor $1-\theta$ can be absorbed in the upper self-test
bound. A uniform conclusion needs $\theta$ bounded away from $1$ and
$o(V)/V\to0$.

Test a proposed source budget on closed K\"ahler--Einstein manifolds,
hyperbolic products, the actual sheared constructions, and local curvature
spikes. A budget that excludes a known exact Einstein control has failed
either in its hypotheses or in the proposed estimate.

\subsection{The same-object intersection of rigidity and flexibility}\label{sub:experiment-C}
For the closed sheared manifolds, $E_D$, negative-orbifold resolutions,
and critical surgery families, compare on each object the quantities
$T$, actual special classes, all ordinary classes, capacities,
$e_\sigma$, and restrictions on exact Einstein outputs.
For a proposed surgery changing $T/\alpha^2$, the input metric change,
extension of actual monopoles, and exclusion or survival of every
Einstein output must be proved together. If these demands conflict,
identify the first incompatible formula and formulate a surgery-cost
or class-failure statement. A different success could combine a
non-monopole global obstruction with the existing ratio-one
constructions. The integral theory then explains why it permits the
approximate input; it need not force every counterexample into its own
critical band.

\section{Literature and comparison of hypotheses}\label{sec:literature}
LeBrun's 2026 survey \cite{LeBrunSurvey} provides a current map of
Einstein four-manifolds. The curvature functionals, actual-monopole
rescaling, convex geometry, and connected-sum inputs used here come from
\cite{LeBrun01,LeBrunConvex,IL,Ishida,Kotschick}. Comparisons should
identify the object, actual class hypothesis, error norm, constant
dependence, and attainment question.

Baldauf's weighted monopoles and harmonic spinors in Ricci flow
\cite{Baldauf} suggest comparisons between amplitude weights and dynamic
entropy. Those equations and conclusions differ from the fixed-metric
clipping and self-test here. The regularity tools in \cite{CT,CN,JN}
continue to apply to controlled regions and exact limits. Huang--Ozuch
\cite{HO} prove isolation results for spherical and hyperbolic Einstein
four-orbifolds as an input to five-dimensional regularity; that is not an
isolation theorem for arbitrary metrics with small $L^2$ Ricci error.
The author is Yiqi Huang. The cited Ishida paper is \emph{Einstein metrics
and exotic smooth structures}; its title is corrected here.

No bibliographic search or theorem count determines priority. A proposed
extension needs comparison at the level of precise propositions, and
an elementary reorganization such as Proposition~\ref{prop:self-test}
should not be confused with supplying its geometric hypotheses.

\section{Further outcomes and the full problem}\label{sec:outcomes}
The integral near-equality theorem, its selected-class cost, and its
special-capacity form provide a complete starting argument. Further
mathematical outcomes may enlarge the actual class or topology range,
produce a realized countermodel to an attempted extension, establish a
source or boundary self-test for a different equation, or connect a
construction or existence argument to a sign branch of
\eqref{eq:branches}. None of these outcomes prescribes a fixed priority
or requires retaining the current proof architecture.

For any proposed object, determine whether its actual classes and square
bounds satisfy the existing theorem. If so, its prohibited approximate
input is already excluded. If not, identify the missing hypothesis and
investigate selection cost, transition cost, another equation, or direct
existence. Exact and degenerating control examples are useful for every
new mechanism. The full Yau26 objective remains open in this article;
the integral obstruction is a proved component of the supplied research,
with an explicitly delimited scope.

\appendix
\section{Source arguments and inheritance}\label{sec:sources}
The standalone argument above integrates the September 21 rigidity
manuscript and the connected English synthesis of the twelve rounds.
The original round files remain unchanged under
\texttt{../references/rounds/}; the September 16 baseline remains in this
research directory. The following source map retains independent proofs,
qualified results, reviews, and open mechanisms. Source scripts verify
finite algebra, coefficients, and examples; they do not replace PDE
existence, nonlinear remainders, or global topology.

\small
\begin{longtable}{L{.35\textwidth}L{.59\textwidth}}
\toprule Source & Content and scope\\\midrule\endhead
\baseref{20260916_rigidity.tex}{20260916\_rigidity.tex};
\baseref{20260916_uniform_gap.tex}{20260916\_uniform\_gap.tex};
\baseref{20260916.md}{20260916.md} & Exact gap, approximate conditional
routes, local flow calculations, and original quantifiers. Retained
independently of the selected-class integral theorem.\\
\roundref{rigidity.md}{rigidity.md};
\roundref{rigidity_review.md}{rigidity\_review.md} & Local spectral
polynomial, scalar multiplier obstruction, Cotton counterexample, cubic
concentration, and norm/first-jet review.\\
\roundref{fixed_ratio.md}{fixed\_ratio.md};
\roundref{curvature_box.md}{curvature\_box.md} & Exact finite-dimensional
Ricci-box optimizations; connector consequences require matched budgets.\\
\roundref{round2_rigidity.md}{round2\_rigidity.md};
\roundref{round2_rigidity_review.md}{round2\_rigidity\_review.md} & Harmonic
and eigenline localization, Walsh phase averaging, spinor model, endpoint
failure, surviving energy, full-curvature-bounded geometric gap.\\
\roundref{round2_weak_curvature.md}{round2\_weak\_curvature.md};
\roundref{round2_weak_review.md}{round2\_weak\_review.md} & Injective Weyl
symbol, first-jet coercivity, weak Cotton, controlled-chart transfer,
buffered developments and multiplicity. No arbitrary-collapse regularity.\\
\roundref{round3_monopole_input.md}{round3\_monopole\_input.md} & Conformal
rescaling, four defects, moments, lowest-eigenvalue inequality, quadratic
spinor differentiation, weighted Hodge and first-jet identities.\\
\roundref{round3_rigidity.md}{round3\_rigidity.md};
\roundref{round3_joint_gap.md}{round3\_joint\_gap.md};
\roundref{round3_general_monopole.md}{round3\_general\_monopole.md} &
Intrinsic Hessian closure, blow-up selection, and controlled-square
extension; pointwise Ricci hypotheses subsequently weakened.\\
\roundref{round3_rigidity_review.md}{round3\_rigidity\_review.md};
\roundref{round3_rigidity_second_review.md}{round3\_rigidity\_second\_review.md}
& Hessian and nonabelian commutator audits, curvature representation,
smoothing order, determinant budget, legitimate self-test.\\
\roundref{round4_rigidity.md}{round4\_rigidity.md};
\roundref{round4_all_sign_gap.md}{round4\_all\_sign\_gap.md} & Main integral
proof, clipped weight, actual/ideal moments, topology transfer, and
all-sign optimization.\\
\roundref{round4_rigidity_review_global.md}{round4\_rigidity\_review\_global.md};
\roundref{round4_global_all_sign_review.md}{round4\_global\_all\_sign\_review.md}
& Unrestricted scalar signs, constants and alternative-square branch;
only the latter file's all-sign part enters this proof.\\
\roundref{round5_rigidity.md}{round5\_rigidity.md} & Ordinary/generic/special
distinction, minimax, convex countermodel, connected sums, ratio-one
bundles, and hyperbolic $A_1$ accounting.\\
\roundref{round5_construction_rigidity_review.md}{\nolinkurl{round5_construction_rigidity_review.md}}
& Receiver tightness and detector convergence; construction-side review,
not an input to the integral theorem.\\
\roundref{round7_endpoint.md}{round7\_endpoint.md} & Filling-independent
concentration identity, bulk moment hypotheses, symmetric bulks,
opposite K\"ahler slope, and modulus-versus-spectrum distinction.\\
\roundref{round7_rigidity.md}{round7\_rigidity.md} & Conditional
nonabelian reduction and literal projection algebra; identification with
all published conventions remains unresolved.\\
\roundref{round8_nonabelian.md}{round8\_nonabelian.md};
\roundref{round9_nonabelian.md}{round9\_nonabelian.md};
\roundref{round10_nonabelian.md}{round10\_nonabelian.md} & Line curvature
cost, partial Dirac/trace seeds, and full-rank residual limitation;
none establishes existence of the full coupled system.\\
\roundref{round11_rankone.md}{round11\_rankone.md};
\roundref{round11_variational.md}{round11\_variational.md} & Rank-one
mixing and second fundamental form; exact variational infima with
distinct attainment questions.\\
\roundref{round12_global.md}{round12\_global.md};
\roundref{round12_uniformization.md}{round12\_uniformization.md};
\roundref{round12_control.md}{round12\_control.md} & Basic classes and
capacities, corrected central block connection and graded energy,
conditional classification, and Schoen-surface control.\\
\bottomrule
\end{longtable}
\normalsize
The paired \texttt{*check*.py} and \texttt{*verification.md} files preserve
the coefficient and finite-dimensional checks with their source versions.
The historical Fable material separately distinguishes the valid OT
topology extension, chirality applications, negative-sign GH separation,
and independent coordinate calculations from the failed problem-W
patching and its specified-cutoff $1/3$ leading term. None is promoted
to an unrestricted construction or a proof of the present main theorem.

\begin{thebibliography}{99}
\bibitem{LeBrun01} C.~LeBrun, \emph{Ricci Curvature, Minimal Volumes, and
Seiberg--Witten Theory}, Invent. Math. 145 (2001), 279--316.
\url{https://arxiv.org/abs/math/0003068}.
\bibitem{LeBrunConvex} C.~LeBrun, \emph{Four-Manifolds, Curvature Bounds,
and Convex Geometry}. \url{https://arxiv.org/abs/math/0611450}.
\bibitem{LeBrunSurvey} C.~LeBrun, \emph{An Overview of Einstein
4-Manifolds}, preprint, 2 September 2026.
\url{https://arxiv.org/abs/2609.03033}.
\bibitem{Ishida} M.~Ishida, \emph{Einstein metrics and exotic smooth
structures}, Pacific J. Math. 258 (2012), 327--349.
\url{https://msp.org/pjm/2012/258-2/pjm-v258-n2-p05-s.pdf}.
\bibitem{IL} M.~Ishida and C.~LeBrun, \emph{Curvature, Connected Sums,
and Seiberg--Witten Theory}. \url{https://arxiv.org/abs/math/0111228}.
\bibitem{Kotschick} D.~Kotschick, \emph{Monopole classes and Einstein
metrics}. \url{https://arxiv.org/abs/math/0306022}.
\bibitem{CT} J.~Cheeger and G.~Tian, \emph{Curvature and injectivity
radius estimates for Einstein 4-manifolds}, J. Amer. Math. Soc. 19 (2006).
\url{https://doi.org/10.1090/S0894-0347-05-00511-4}.
\bibitem{CN} J.~Cheeger and A.~Naber, \emph{Regularity of Einstein
Manifolds and the Codimension 4 Conjecture}.
\url{https://arxiv.org/abs/1406.6534}.
\bibitem{JN} W.~Jiang and A.~Naber, \emph{$L^2$ curvature bounds on
manifolds with bounded Ricci curvature}, Ann. of Math. 193 (2021).
\url{https://annals.math.princeton.edu/2021/193-1/p02}.
\bibitem{Derd} A.~Derdzi\'nski, \emph{Self-dual K\"ahler manifolds and
Einstein manifolds of dimension four}, Compositio Math. 49 (1983), 405--433.
\url{https://www.numdam.org/item/CM_1983__49_3_405_0/}.
\bibitem{Baldauf} J.~Baldauf, \emph{Harmonic spinors in the Ricci flow},
J. Geom. Anal. (2024).
\url{https://arxiv.org/abs/2210.14198};
\url{https://doi.org/10.1007/s12220-024-01665-y}.
\bibitem{HO} Yiqi Huang and T.~Ozuch, \emph{Regularity of Einstein
5-manifolds via 4-dimensional gap theorems}, version 2, 16 February 2026.
\url{https://arxiv.org/abs/2512.21317v2}.
\end{thebibliography}
\end{document}
